% APPEND BLANK -- os/samir/cmd/mutate.c m_append (BLANK required; adds a
% blank record, makes it current, inserts its key into every open index).
% Gate: test-interp-replace, test-samir-write.
% APPEND FROM -- os/samir/cmd/mutate.c d_append_from (Using III+ U5-28/PDF
% 140: fields matched by name in any order; deleted records skipped only
% under SET DELETED ON; default .DBF; the database in USE refused, 3; a
% memo field in the database in USE, 34; a field of the same name and
% another type, 9; every open index rebuilt).  Lane samirdml 14a8fc75.
% Gates: test-samir-dml, test-samir-dml-emu.
\begin{XbaseEntry}{APPEND}{command}
Adds records to the end of the database in USE: one blank record, or the
records of another database.
\Syntax
\SyntaxLine{\kw{APPEND BLANK}}
\par or\par
\SyntaxLine{\kw{APPEND FROM} \xp{file}}
\Defaults
\kw{APPEND FROM} assumes the extension \file{.DBF}.
\Usage
\cmd{APPEND BLANK} adds one record to the end of the database in the selected
work area and makes it the current record.  Every field of the new record is
empty: character fields are blanks, numeric fields are zero, date fields are
blank and logical fields are false.  The record count, \fn{RECCOUNT()}, grows
by one.

The new record is written to the disk at once.  If index files are open with
the database, the new record is entered in each of them.

Fill in the fields of the new record with \cmd{REPLACE}.  In a program, an
\cmd{ACCEPT} or \cmd{INPUT} for each field followed by \cmd{APPEND BLANK} and
\cmd{REPLACE} enters a record.

\cmd{APPEND FROM} adds every record of the database named to the end of the
database in USE, in the order of its records.  Fields are matched by name:
a field of the database in USE takes the value of the field of the same name
in the other database, wherever it stands in its structure.  A field the
other database does not have is left empty, and a field the database in USE
does not have is not copied.  A character value longer than its field is cut
off on the right.  Records marked for deletion are appended, and are not
marked in the database in USE, unless SET DELETED is ON: then they are left
out.  The last record appended becomes the current record, and every open
index file is brought up to date.

The database named must not be the database in USE
(\xmsg{3}{File is already open.}), and two fields of the same name must have
the same type (\xmsg{9}{Data type mismatch.}).  Records cannot be appended to
a database that has a memo field (\xmsg{34}{Operation with Memo field
invalid.}).
\Examples
A new record is added and two of its fields are filled in:
\DotExample{examples/c-append}
The batch of new invoices on the sample diskette, \file{ARBATCH.DBF}, has the
fields of the ledger in another order, and no \texttt{PAID} field.  Its three
records are added to \file{INVOICE.DBF}:
\DotExample{examples/c-append-2}
\SeeAlso{REPLACE, RECCOUNT(), DELETE, COPY FILE}
\end{XbaseEntry}

